\writeSectionFrame{\presentationTitle}{Fluent API}
\begin{frame}{Was ist eine Fluent API?}
	\begin{itemize}
		\item Eine Programmierschnittstelle mit flüssiger, lesbarer Syntax
		\item Methoden liefern meist das Objekt selbst oder ein Folgeobjekt zurück
		\item Ermöglicht Methodenverkettung (method chaining)
		\item Fokus liegt auf Lesbarkeit und Ausdrucksstärke
	\end{itemize}
\end{frame}

\begin{frame}{Motivation}
	\begin{itemize}
		\item Klassische APIs sind oft:
		\begin{itemize}
			\item schwer lesbar
			\item technisch, aber nicht semantisch
		\end{itemize}
		\item Fluent APIs sollen sich eher wie eine Domänensprache anfühlen
		\item Code soll sich lesen lassen wie ein Satz
	\end{itemize}
\end{frame}

\begin{frame}[fragile]{Klassische API}
	\begin{lstlisting}[language=Java]
Query q = new Query();
q.setTable("users");
q.addCondition("age > 18");
q.setLimit(10);
	\end{lstlisting}
	
	\pause
	\begin{itemize}
		\item Technisch korrekt
		\item Reihenfolge und Bedeutung nur implizit
	\end{itemize}
\end{frame}

\begin{frame}[fragile]{Fluent API}
	\begin{lstlisting}[language=Java]
Query q = Query
	.from("users")
	.where("age > 18")
	.limit(10);
	\end{lstlisting}
	
	\pause
	\begin{itemize}
		\item Lesbar von oben nach unten
		\item Reihenfolge wird Teil der Bedeutung
		\item Intention klar erkennbar
	\end{itemize}
\end{frame}

\begin{frame}{Technische Grundlage}
	\begin{itemize}
		\item Methoden geben:
		\begin{itemize}
			\item \texttt{this} zurück oder
			\item ein anderes Objekt derselben API
		\end{itemize}
		\item Kein neues Sprachfeature
		\item Nutzt:
		\begin{itemize}
			\item Rückgabewerte
			\item Interfaces
			\item Immutable Objekte (häufig)
		\end{itemize}
	\end{itemize}
\end{frame}

\begin{frame}{Typische Einsatzgebiete}
	\begin{itemize}
		\item Builder (\texttt{StringBuilder}, \texttt{HttpRequest})
		\item Datenbankabfragen (JPA Criteria, jOOQ)
		\item Konfigurationen
		\item Test-Frameworks (z.B. AssertJ oder ArchUnit)
	\end{itemize}
\end{frame}
